\scanpage{8}\mspage{\handwritten{23}}
\makeatletter\tagsleft@true\makeatother
\[
  0 \handwritten{\longrightarrow} \underline{\mathrm Z}/\handwritten{\delta''}\,\underline{\mathrm Z}
  \handwritten{\longrightarrow \mathord{\tikz[x=1em,y=1ex,baseline=0pt,line width=.045em,line cap=round]{\draw (0,.08)--(.88,.08) (0,1.43)--(.88,1.43) (.19,.08)--(.19,1.43) (.44,.08)--(.44,1.43) (.69,.08)--(.69,1.43);}}'}(Y,\underline A)
  \handwritten{\longrightarrow \mathord{\tikz[x=1em,y=1ex,baseline=0pt,line width=.045em,line cap=round]{\draw (0,.08)--(.88,.08) (0,1.43)--(.88,1.43) (.19,.08)--(.19,1.43) (.44,.08)--(.44,1.43) (.69,.08)--(.69,1.43);\draw[line width=.02em] (.44,.73)--(.69,.73);}}^{\deleted{\unsure{\prime|*}}}}(Y,\underline B)
  \handwritten{\longrightarrow} 0 \quad . \tag*{(4.38)}
\]

Explicitons deux cas particuliers importants:

\deleted{(\unsure{4.6})}
\uline{Cas particulier} \tikz[baseline=0pt,x=1em,y=1ex,line cap=round]{
  \path[use as bounding box] (0,-.7) rectangle (.15,4.3);
  \node[overlay,anchor=base west,inner sep=0pt] at (.05,2.7) {\inserted{(4.6):}};
  \draw[overlay] (0,-.7) -- (0,1.9) .. controls (0,2.3) and (.2,2.3) .. (.5,2.3)
    -- (4.6,2.3) .. controls (4.9,2.3) and (5.1,2.5) .. (5.1,2.9);
}\deleted{1} $Y$ \uline{est une courbe algébrique sur un corps algébriquement clos}. Alors la \deleted{relation (4.13) est automatiqu\unsure{e.|e. La}} plus grande partie de la discussion précédente devient triviale. La relation (4.13) est automatiquement satisfaite, et on a (4.12 bis) :
\[
  \mathrm{Br}(X)\handwritten{\overset{\sim}{\longrightarrow}}H^{1}(Y,\underline B)\quad ; \tag*{(4.39)}
\]
lorsque de plus les $\handwritten{\delta}_{y}$ sont égaux à $1$, i.e. \handwritten{si} $X$ n'a pas de fibre multiple sur $Y$, alors le deuxième membre s'explicite \deleted{comme étant égal à} par
\[
  0\handwritten{\longrightarrow}\underline{\mathrm Z}/\handwritten{\delta}\underline{\mathrm Z}
  \handwritten{\longrightarrow}H^{1}(Y,\underline{\deleted{\unsure{X|B}}\handwritten{A}})
  \handwritten{\longrightarrow}H^{1}(Y,\underline B)\handwritten{\longrightarrow}0\quad , \tag*{(4.40)}
\]
compte tenu que \deleted{dans} l'on a ici nécessairement $\handwritten{\delta}=\handwritten{\delta'}=\handwritten{\delta''}$ .

\uline{Cas particulier} (4.7): $Y$ \uline{est une courbe algébrique propre sur un corps fini} $k$. Ici encore, (4.13) \tikz[baseline=0pt,x=1em,y=1ex,line cap=round]{
  \path[use as bounding box] (0,-.8) rectangle (.1,4.6);
  \node[overlay,anchor=base west,inner sep=0pt] at (.15,2.8) {\deleted{\handwritten{\unsure{et (4.16)}}}};
  \draw[overlay] (.15,-.7) .. controls (-.4,1.8) and (0,2.0) .. (1,2.0)
    -- (6.5,2.0) .. controls (6.8,2.0) and (6.9,2.2) .. (6.8,2.7);
}\deleted{\unsure{sont|est}}\handwritten{est} automatiquement satisfaite, $k$ étant parfait, \deleted{et}\inserted{\handwritten{ainsi que (4.16) puisque $k(y)$ est de dimension $\leq 1$, d'où $H^{3}(k(y),\mathbb{G}_{\mathrm m})=0$, de plus}} on aura $\handwritten{\delta}=\handwritten{\delta'}\handwritten{\deleted{\unsure{=\delta''}}}$, et de même pour les invariants hensélisés et strictement hensélisés. Comme on a encore $\mathrm{Br}(Y)=0$, on trouve ici (4.4):
\[
  0\handwritten{\longrightarrow}\mathrm{Br}(X)\handwritten{\longrightarrow \mathord{\tikz[x=1em,y=1ex,baseline=0pt,line width=.045em,line cap=round]{\draw (0,.08)--(.88,.08) (0,1.43)--(.88,1.43) (.19,.08)--(.19,1.43) (.44,.08)--(.44,1.43) (.69,.08)--(.69,1.43);}}}(Y,\underline B)
  \handwritten{\longrightarrow}\underline{\deleted{\unsure{2|Z}}\handwritten{\mathrm Z}}/\deleted{\unsure{1|\delta}}\handwritten{\Delta}\underline{\mathrm Z}
  \handwritten{\longrightarrow}0\quad ,\qquad\text{où}\quad\handwritten{\Delta\mid\delta}\quad . \tag*{(4.41)}
\]
\deleted{\unsure{compte tenu de|compte tenu du}} D'autre part, si les $\handwritten{\delta}_{y}$ sont égaux à $1$, i.e. si le morphisme déduit de $f:X\handwritten{\longrightarrow}Y$ en passant à la clôture algébrique de $k$ n'a pas de fibre multiple, alors on trouve la suite exacte
\[
  0\handwritten{\longrightarrow}\underline{\mathrm Z}/\handwritten{\delta''}\underline{\mathrm Z}
  \handwritten{\longrightarrow \mathord{\tikz[x=1em,y=1ex,baseline=0pt,line width=.045em,line cap=round]{\draw (0,.08)--(.88,.08) (0,1.43)--(.88,1.43) (.19,.08)--(.19,1.43) (.44,.08)--(.44,1.43) (.69,.08)--(.69,1.43);}}}(Y,\underline A)
  \handwritten{\to \mathord{\tikz[x=1em,y=1ex,baseline=0pt,line width=.045em,line cap=round]{\draw (0,.08)--(.88,.08) (0,1.43)--(.88,1.43) (.19,.08)--(.19,1.43) (.44,.08)--(.44,1.43) (.69,.08)--(.69,1.43);}}}(Y,\underline B)\handwritten{\longrightarrow}0\quad , \tag*{(4.42)}
\]
ce qui est un cas particulier de (4.37), compte tenu de\deleted{s} \tikz[baseline=0pt,x=1em,y=1ex,line cap=round]{
  \path[use as bounding box] (0,-.5) rectangle (.1,4.0);
  \node[overlay,anchor=base west,inner sep=0pt] at (.1,2.4) {\inserted{\handwritten{la}}};
  \draw[overlay] (0,-.5) .. controls (.4,.8) and (.6,1.6) .. (1.1,1.7)
    .. controls (1.5,1.8) and (1.9,2.0) .. (2.0,2.4);
} relation\deleted{s} suivante\deleted{s}\unsure{:}
\[
  \handwritten{\bar{\delta}}_{y}=\handwritten{\delta}_{y} \tag*{(4.43)}
\]
pour tout point fermé $y$ de $Y$, qui implique ici $\handwritten{\bar{\delta}}_{y}=1$ d'où $\handwritten{\mathord{\tikz[x=1em,y=1ex,baseline=0pt,line width=.045em,line cap=round]{\draw (0,.08)--(.88,.08) (0,1.43)--(.88,1.43) (.19,.08)--(.19,1.43) (.44,.08)--(.44,1.43) (.69,.08)--(.69,1.43);}}'}(Y,\underline A)=\handwritten{\mathord{\tikz[x=1em,y=1ex,baseline=0pt,line width=.045em,line cap=round]{\draw (0,.08)--(.88,.08) (0,1.43)--(.88,1.43) (.19,.08)--(.19,1.43) (.44,.08)--(.44,1.43) (.69,.08)--(.69,1.43);}}}(Y,\underline A)$.

La démonstration de\deleted{s} \handwritten{l'}égalité\deleted{s} (4.43), qui est une conséquence facile du
\reversemarginpar\margin{\handwritten{\unsure{x|$\times$}}}\normalmarginpar
théorème de \deleted{Lang}\handwritten{LANG} \tikz[baseline=0pt,x=1em,y=1ex,line cap=round]{
  \path[use as bounding box] (0,-.6) rectangle (.15,4.5);
  \node[overlay,anchor=base west,inner sep=0pt] at (-5.0,2.7) {\inserted{\handwritten{[24] $H^{1}(k,G)=0$}}};
  \draw[overlay] (-5.5,2.6) .. controls (-6.0,1.6) and (-5.7,1.5) .. (-4.7,1.5)
    -- (-.55,1.2) .. controls (-.2,1.2) and (-.05,1.2) .. (0,.65)
    .. controls (.15,1.3) and (.4,1.4) .. (.7,1.4)
    -- (4.8,1.5) .. controls (5.1,1.5) and (5.2,1.5) .. (5.2,2.6);
  \draw[overlay] (0,.65) -- (0,-.6);
} pour le corps fini $k(y)$ et le schéma de Picard \deleted{\unsure{con|Con}}\handwritten{con}nexe $\handwritten{G=}\underline{\mathrm{Pic}}^{0}_{X_{y}/y}$, est laissé\handwritten{e} au lecteur comme un exercice plaisant et délectable.
% archive sheet number: 7
\clearpage
